\documentclass{beamer}
\usepackage{tikz,amsmath,hyperref,graphicx,stackrel,animate,amssymb}
\usetikzlibrary{positioning,shadows,arrows,shapes,calc}
\newcommand{\argmax}{\operatornamewithlimits{argmax}}
\newcommand{\argmin}{\operatornamewithlimits{argmin}}
\mode<presentation>{\usetheme{Frankfurt}}
\AtBeginSection[]
{
  \begin{frame}<beamer>
    \frametitle{Outline}
    \tableofcontents[currentsection,currentsubsection]
  \end{frame}
}
\title{Lecture 12: Circular Convolution}
\author{Mark Hasegawa-Johnson}
\date{ECE 401: Signal and Image Analysis}
\begin{document}

% Title
\begin{frame}
  \maketitle
\end{frame}

% Title
\begin{frame}
  \tableofcontents
\end{frame}

%%%%%%%%%%%%%%%%%%%%%%%%%%%%%%%%%%%%%%%%%%%%
\section[Review]{Review: DFT}
\setcounter{subsection}{1}

\begin{frame}
  \frametitle{Review: DFT}

  The DFT (discrete Fourier transform) of any signal is
  $X[k]$, given by
  \begin{align*}
    X[k] &= \sum_{n=0}^{N-1} x[n]e^{-j\frac{2\pi kn}{N}}\\
    x[n] &= \frac{1}{N}\sum_0^{N-1} X[k]e^{j\frac{2\pi kn}{N}}
  \end{align*}
  Particular useful examples include:
  \begin{align*}
    f[n]=\delta[n] &\leftrightarrow F[k]=1\\
    g[n]=\delta\left[(\!(n-n_0)\!)_N\right] &\leftrightarrow G[k]=e^{-j\frac{2\pi kn_0}{N}}
  \end{align*}
\end{frame}

\begin{frame}
  \frametitle{Properties of the DTFT}

  Properties worth knowing  include:
  \begin{enumerate}
    \setcounter{enumi}{-1}
  \item Periodicity: $X[k+N]=X[k]$
  \item Conjugate symmetric: $x[n]$ Real $\leftrightarrow X[k]=X^*[-k]$
  \item Linearity:
    \[z[n]=ax[n]+by[n]\leftrightarrow Z[k]=aX[k]+bY[k]
    \]
  \item Time Shift: $x[n-n_0]\leftrightarrow e^{-j\frac{2\pi
      kn_0}{N}}X[k]$, paying attention to the fact that $x[n]$ is
    periodic in time
  \item Frequency Shift: $e^{j\frac{2\pi k_0 n}{N}}x[n]\leftrightarrow
    X[k-k_0]$, paying attention to the fact that $X[k]$ is periodic
    in frequency
  \end{enumerate}
\end{frame}

\begin{frame}
  \frametitle{Convolution}

  Convolution (finite impulse response filtering) is a generalization of weighted local averaging:
  \begin{displaymath}
    y[n] = h[n]\ast x[n] \equiv \sum_m x[m] h[n-m] = \sum_m x[n-m] h[m]
  \end{displaymath}
  \begin{itemize}
  \item If all samples of $h[n]$ are positive, then it's a weighted
    local averaging filter
  \item If the samples of $h[n]$ are positive for $n>0$ and negative
    for $n<0$ (or vice versa), then it's a weighted local differencing
    filter
  \end{itemize}
\end{frame}

\begin{frame}
  \frametitle{Frequency Response}
  \begin{itemize}
  \item {\bf Tones in $\rightarrow$ Tones out}
    \begin{align*}
      x[n]=e^{j\omega n} &\rightarrow y[n]=H(\omega)e^{j\omega n}\\
      x[n]=\cos\left(\omega n\right)
      &\rightarrow y[n]=|H(\omega)|\cos\left(\omega n+\angle H(\omega)\right)\\
      x[n]=A\cos\left(\omega n+\theta\right)
      &\rightarrow y[n]=A|H(\omega)|\cos\left(\omega n+\theta+\angle H(\omega)\right)
    \end{align*}
  \item where the {\bf Frequency Response} is given by
    \[
    H(\omega) = \sum_m h[m]e^{-j\omega m}
    \]
  \end{itemize}
\end{frame}  

\begin{frame}
  \frametitle{Putting it all together}

  If the input signal is
  \begin{align*}
    x[n] = \frac{1}{N}\sum_{n=0}^{N-1} X[k]e^{j\frac{2\pi kn}{N}}
  \end{align*}
  \ldots then the output of $y[n]=h[n]\ast x[n]$ is\ldots
  \begin{align*}
    y[n] &= \frac{1}{N}\sum_{n=0}^{N-1} H[k] X[k]e^{j\frac{2\pi kn}{N}}\\
    H[k] &\equiv H\left(\frac{2\pi k}{N}\right)
  \end{align*}
\end{frame}

%%%%%%%%%%%%%%%%%%%%%%%%%%%%%%%%%%%%%%%%%%%%
\section[FFT]{Computational Complexity: Fast Fourier Transform}
\setcounter{subsection}{1}

\begin{frame}
  \frametitle{Convolution is an $\mathcal{O}\left\{N^2\right\}$ Operation}

  Suppose we convolve an N-sample $h[n]$ with an N-sample $x[n]$.
  Consider carefully where the nonzero terms are:
  \begin{align*}
    y[n] = \sum_{m=n}^{N-1} h[m]x[n-m],~~~0\le n\le 2N-1
  \end{align*}
  Let's do a computational complexity analysis:
  \begin{itemize}
  \item There are $2N=\mathcal{O}\left\{N\right\}$ output samples.
  \item The $m^{\text{th}}$ output sample results from a sum of length
    $N-m=\mathcal{O}\left\{N\right\}$
  \item The total complexity is therefore $\mathcal{O}\left\{N^2\right\}$
  \end{itemize}
\end{frame}

\begin{frame}
  \frametitle{Fast Fourier Transform (FFT)}

  \begin{itemize}
  \item The regular DFT formula, $X[k]=\sum
    x[n]e^{-j\frac{2\pi{kn}}{N}}$, is also
    $\mathcal{O}\left\{N^2\right\}$ (N terms in the sum, N output coefficients).
  \item In 1965, Cooley and Tukey found a way to re-arrange the terms
    of the DFT so it can be computed in
    $\mathcal{O}\left\{N\log_2{N}\right\}$ multiply-add operations.
    Their algorithm is called the ``fast Fourier transform'' (FFT).
  \item For example, suppose $N=1024$, then
    \begin{itemize}
    \item $N^2=1,048,576$ operations
    \item $N\log_2 N=10,240$ operations (a factor of 100 savings)
    \end{itemize}
  \end{itemize}
\end{frame}

\begin{frame}
  \frametitle{Can we use the FFT to speed up convolution?}

  Remember this formula:
  \begin{align*}
    y[n] &= \frac{1}{N}\sum_{n=0}^{N-1} H[k] X[k]e^{j\frac{2\pi kn}{N}}
  \end{align*}
  How about the following?
  \begin{enumerate}
  \item Take the FFT of $x[n]$ to find $X[k]$ ($\mathcal{O}\left\{N\log N\right\}$)
  \item Take the FFT of $h[n]$ to find $H[k]$ ($\mathcal{O}\left\{N\log N\right\}$)
  \item Multiply them to find $Y[k]$ ($\mathcal{O}\left\{N\right\}$)
  \item Take the inverse FFT to find $y[n]$ ($\mathcal{O}\left\{N\log N\right\}$)
  \end{enumerate}
  Is the total complexity $\mathcal{O}\left\{N\log N\right\}$, much
  faster than regular convolution?  YES!!  Does it compute the desired
  answer? NO!!! Why not, and how can we fix it?
\end{frame}


%%%%%%%%%%%%%%%%%%%%%%%%%%%%%%%%%%%%%%%%%%%%
\section[Circular Convolution]{Circular Convolution}
\setcounter{subsection}{1}

\begin{frame}
  \frametitle{Multiplying two DFTs: what we think we're doing}

  \centerline{\includegraphics[width=\textwidth]{exp/convolutionlinear.png}}
\end{frame}

\begin{frame}
  \frametitle{Multiplying two DFTs: what we're actually doing}

  \centerline{\includegraphics[width=\textwidth]{exp/convolutionperiodic.png}}
\end{frame}

\begin{frame}
  \frametitle{Circular convolution}

  Suppose $Y[k]=H[k]X[k]$, then
  \begin{align*}
    y[n] &= \frac{1}{N}\sum_{k=0}^{N-1} H[k]X[k]e^{j\frac{2\pi kn}{N}}\\
    &= \frac{1}{N}\sum_{k=0}^{N-1} H[k]\left(\sum_{m=0}^{N-1} x[m]e^{-j\frac{2\pi km}{N}}\right)
    e^{j\frac{2\pi kn}{N}}\\
    &= \sum_{m=0}^{N-1}x[m]\left(\frac{1}{N}\sum_{k=0}^{N-1} H[k]e^{-j\frac{2\pi k(n-m)}{N}}\right)\\
    &= \sum_{m=0}^{N-1} x[m] h\left[(\!(n-m)\!)_N\right]
  \end{align*}
  \ldots where the notation $(\!(n-m)\!)_N$ means ``modulo $N$,'' i.e., we need to pay attention to the periodicity.
\end{frame}

\begin{frame}
  \frametitle{Circular convolution}

  Multiplying the DFT means {\bf circular convolution} of the time-domain signals:
  \begin{displaymath}
    y[n]=h[n]\circledast x[n] \leftrightarrow Y[k] = H[k]X[k],
  \end{displaymath}
  
  Circular convolution ($h[n]\circledast x[n]$) is defined like this:
  \begin{displaymath}
    h[n]\circledast x[n] = \sum_{m=0}^{N-1}x[m]h\left[(\!(n-m)\!)_N\right]
    = \sum_{m=0}^{N-1}h[m]x\left[(\!(n-m)\!)_N\right]
  \end{displaymath}
\end{frame}

\begin{frame}
  \frametitle{Linear convolution example}

  \centerline{\animategraphics[loop,controls,width=\textwidth]{10}{exp/linear}{0}{68}}
\end{frame}

\begin{frame}
  \frametitle{Circular convolution example}

  \centerline{\animategraphics[loop,controls,width=\textwidth]{10}{exp/circ}{0}{53}}
\end{frame}


%\begin{frame}
%  \frametitle{Circular convolution example}

%  \centerline{\animategraphics[loop,controls,width=\textwidth]{10}{exp/circularsquares}{0}{49}}
%\end{frame}

%\begin{frame}
%  \frametitle{Circular convolution example}

%  \centerline{\animategraphics[loop,controls,width=\textwidth]{10}{exp/circularexps}{0}{49}}
%\end{frame}

\begin{frame}
  \frametitle{Quiz!}

  Go to the course webpage, try the quiz!
\end{frame}


%%%%%%%%%%%%%%%%%%%%%%%%%%%%%%%%%%%%%%%%%%%%
\section[Zero-Padding]{Zero-Padding}
\setcounter{subsection}{1}

\begin{frame}
  \frametitle{How long is $h[n]\ast x[n]$?}

  If $x[n]$ is $M$ samples long, and $h[n]$ is $L$ samples long, then
  their linear convolution, $y[n] = x[n]\ast h[n]$, is $M+L-1$ samples long.
  
  \centerline{\includegraphics[width=\textwidth]{exp/convolutionlengths.png}}
\end{frame}

\begin{frame}
  \frametitle{Zero-padding turns circular convolution into linear convolution}

  How it works:
  \begin{itemize}
  \item $h[n]$ is length-$L$
  \item $x[n]$ is length-$M$
  \item As long as they are both zero-padded to length $N\ge L+M-1$, then
  \item $y[n] = h[n]\circledast x[n]$ is the same as $h[n]\ast x[n]$.
  \end{itemize}
\end{frame}
  
\begin{frame}
  \frametitle{Zero-padding turns circular convolution into linear convolution}

  Why it works:  Either\ldots
  \begin{itemize}
  \item $n-m$ is a positive number, between $0$ and $N-1$.  Then
    $(\!(n-m)\!)_N=n-m$, and therefore
    \begin{displaymath}
      x[m]h\left[(\!(n-m)\!)_N\right]=x[m]h[n-m]
    \end{displaymath}
  \item $n-m$ is a negative number, between $0$ and $-(L-1)$.  Then
    $(\!(n-m)\!)_N=N+n-m\ge N-(L-1)>M-1$, so
    \begin{displaymath}
      x[m]h\left[(\!(n-m)\!)_N\right]=0
    \end{displaymath}
  \end{itemize}
\end{frame}
  
\begin{frame}
  \frametitle{Case \#1: $n-m$ is positive, so circular
    convolution is the same as linear convolution}

  \centerline{\includegraphics[width=\textwidth]{exp/circandlinear50.png}}
\end{frame}
  
\begin{frame}
  \frametitle{Case \#2: $n-m$ is negative, so it wraps around, but $N$
    is long enough so that the wrapped part of
    $h\left[(\!(n-m)\!)_N\right]$ doesn't overlap with $x[m]$}
  
  \centerline{\includegraphics[width=\textwidth]{exp/circandlinear0.png}}
\end{frame}
  
\begin{frame}
  \frametitle{Zero-padding turns circular convolution into linear convolution}

  \centerline{\animategraphics[loop,controls,width=\textwidth]{10}{exp/circandlinear}{0}{68}}
\end{frame}



%%%%%%%%%%%%%%%%%%%%%%%%%%%%%%%%%%%%%%%%%%%%
\section[Summary]{Summary}
\setcounter{subsection}{1}

\begin{frame}
  \frametitle{Circular convolution}

  Multiplying the DFTs, instead of implementing convolution directly,
  reduces computation from $\mathcal{O}\left\{N^2\right\}$ to
  $\mathcal{O}\left\{N\log N\right\}$.
  However, multiplying the
  DFT means {\bf circular convolution} of the time-domain signals:
  \begin{displaymath}
    y[n]=h[n]\circledast x[n] \leftrightarrow Y[k] = H[k]X[k],
  \end{displaymath}
  
  Circular convolution ($h[n]\circledast x[n]$) is defined like this:
  \begin{displaymath}
    h[n]\circledast x[n] = \sum_{m=0}^{N-1}x[m]h\left[(\!(n-m)\!)_N\right]
    = \sum_{m=0}^{N-1}h[m]x\left[(\!(n-m)\!)_N\right]
  \end{displaymath}

  Circular convolution is the same as linear convolution if and only if $N\ge L+M-1$.
\end{frame}


\end{document}
